\documentclass[11pt]{article}
\usepackage[a4paper,margin=26mm]{geometry}
\usepackage{amsmath,amssymb,amsthm,mathtools,booktabs,hyperref,xurl}
\setlength{\emergencystretch}{3em}
\newtheorem{theorem}{Theorem}
\newtheorem{lemma}[theorem]{Lemma}
\newtheorem{proposition}[theorem]{Proposition}
\newtheorem{corollary}[theorem]{Corollary}
\newtheorem{remark}[theorem]{Remark}
\newcommand{\Q}{\mathbb Q}
\newcommand{\Z}{\mathbb Z}
\newcommand{\Res}{\operatorname{Res}}
\newcommand{\End}{\operatorname{End}}
\newcommand{\Hom}{\operatorname{Hom}}
\title{Further point-chart reductions and precise obstructions}
\author{\(\mathbb Z\mathbb Q\) Diophantine Project}
\date{23 November 2026}
\begin{document}\maketitle
\begin{abstract}
We record an exact common-point chart reduction and several scoped
obstructions to further identifications. Once the contact gauge is
shared, two private rational coordinates can be removed from each
input chart at the cost of four common coordinates, without changing
any detector or its degree. Separately, the two rational parameters
of the existing elliptic point chart cannot be replaced by one
input-independent rational parameter with finite algebraic fibers:
the relevant restriction of scalars is simple over \(\Q\), and
Mordell--Lang makes every infinite subset of its cyclic orbit
Zariski dense. We also show why the alternate numerator point,
or all of its slope or reciprocal-abscissa data, cannot simply
be identified across distinct input labels. These are statements
about specified constructions, not lower bounds for all
Diophantine representations.
\end{abstract}

\section{An exact reduction of a finite-pair chart}
\label{sec:chart}

Let \(F_0\) be a field of characteristic zero and let \(D\in
F_0^\times\). On
\[
                         E_D:y^2=x^3-Dx
\]
the source finite-pair chart has coordinates \(s,t,F,G,\eta,d\)
and equations
\begin{align}
 d&=F(\eta+s+t),& d&=\eta G,\label{eq:pair1}\\
 D(F+G)&=(s-t)\eta-t^2,&
 1&=d(s-t)-DFG.\label{eq:pair2}
\end{align}
Its point interpretation is
\[
                U=(F^{-1},s/F),\qquad V=(G^{-1},t/G).
\]
The chart includes the tangent case and finite two-torsion
differences; \(d\) is allowed to vanish. The source proof
\cite[Section 3]{physical} gives its full converse.

\begin{lemma}[The common curve identity]\label{lem:identity}
Put
\[
\begin{split}
 r_1&=d-F(\eta+s+t),\\
 r_3&=D(F+G)-(s-t)\eta+t^2,\\
 r_4&=d(s-t)-DFG-1 .
\end{split}
\]
Then the polynomial identity
\begin{equation}\label{eq:identity}
        Fs^2+DF^2-1=r_4-(s-t)r_1+Fr_3
\end{equation}
holds over \(\Z[D,s,t,F,G,\eta,d]\).
\end{lemma}
\begin{proof}
Substitute the three displayed expressions. The terms involving
\(d\), \(FG\), and \(\eta(s-t)F\) cancel. The remaining
quadratic expression in \(s,t\) is
\(F((s-t)(s+t)+t^2)=Fs^2\), giving the identity.
\end{proof}

\begin{proposition}[Eliminating a private chart field]
\label{prop:physical}
Suppose \(k\) copies of the chart share \(s,F\). Introduce
one common field coordinate \(x\), and replace the \(k\)
charts by
\begin{equation}\label{eq:commoncurve}
                    Fx=1,\qquad x=s^2+DF
\end{equation}
and, for \(1\le i\le k\),
\begin{equation}\label{eq:reduced}
 F(\eta_i+s+t_i)=\eta_iG_i,\qquad
 D(F+G_i)=(s-t_i)\eta_i-t_i^2 .
\end{equation}
All these equations are quadratic. Their projection to the
retained coordinates is exactly that of the original shared
charts. The eliminated value is uniquely recovered as
\[
                     d_i=F(\eta_i+s+t_i).
\]
\end{proposition}
\begin{proof}
An original chart satisfies \(Fs^2+DF^2=1\) by
Lemma~\ref{lem:identity}, so \(F\ne0\). Taking \(x=F^{-1}\)
gives \eqref{eq:commoncurve}, while the two formulas for
\(d_i\) give the first equation in \eqref{eq:reduced}.
The second is an unchanged source equation.

Conversely \eqref{eq:commoncurve} gives
\(Fs^2+DF^2=1\). Define \(d_i\) as above. Both equations
in \eqref{eq:pair1} hold, and \(r_3=0\).
Identity \eqref{eq:identity} then gives \(r_4=0\).
Every original chart is recovered. No division by
\(s,t_i,\eta_i,G_i\), or \(d_i\) was introduced. In
particular the zero-\(d_i\) and tangent boundary cases
are preserved.
\end{proof}

For \(F_0=\Q(\sqrt2)\), each field coordinate is represented
by two rational coordinates and each field equation by its
two rational components. Proposition~\ref{prop:physical}
would save \(2k-2\) rational coordinates before the source
Laurent compilation.

\subsection{The obstruction in the old stored scales}

In both quartic interfaces the common source fields are
\[
 s=X_{201}+\sqrt2X_{202},\qquad
 F=X_{205}+\sqrt2X_{206}.
\]
They have contact weight zero. The private fields
\(t,\eta\) also have contact weight zero, whereas
the physical \(G,d\) have weight one:
\[
                          G=c\,\widetilde G,\qquad
                          d=c\,\widetilde d .
\]
All these fields have zero scalar-normalizer weight.
Simply substituting \(G=c\widetilde G\) in the first
equation of \eqref{eq:reduced} produces
\[
              F(\eta+s+t)-c\eta\widetilde G=0.
\]
The last term is cubic in the stored coordinates. Multiplying
the equation by one Laurent monomial cannot make both its
ordinary quadratic terms and its \(c\)-weighted quadratic
term have residual contact exponent zero. Thus direct
deletion of \(d\), with the other scales and the existing
quadratic-unit compilation unchanged, does not establish
a quartic saving.

\subsection{A shared contact gauge removes that obstruction}

\begin{theorem}[Quadratic reduction with a common gauge]
\label{thm:stored}
Suppose the \(k\) input charts have a common nonzero contact
gauge \(c\), common fields \(s,F\), and the private stored
fields \(t_i,\eta_i,\widetilde G_i,\widetilde d_i\) just
described. Assume \(\widetilde d_i\) occurs nowhere outside
its own finite-pair chart.
Introduce two common field coordinates \(x,\widehat F\)
and the three common field equations
\begin{equation}\label{eq:common-scaled}
                  F=c\widehat F,\qquad
                  Fx=1,\qquad x=s^2+DF.
\end{equation}
For each input retain only
\begin{equation}\label{eq:storedreduced}
 \widehat F(\eta_i+s+t_i)=\eta_i\widetilde G_i,\qquad
 D(F+c\widetilde G_i)=(s-t_i)\eta_i-t_i^2 .
\end{equation}
The resulting system is quadratic and has exactly the
same projection as the original system. Over \(\Q(\sqrt2)\)
it saves
\[
                       2k-4\quad\text{rational coordinates}
\]
and \(4k-6\) rational quadratic residuals. For \(k=12\)
the savings are 20 coordinates and 42 residuals.
\end{theorem}
\begin{proof}
Since \(c\ne0\), define
\[
 \widetilde d_i=\widehat F(\eta_i+s+t_i),\qquad
 G_i=c\widetilde G_i,\qquad d_i=c\widetilde d_i .
\]
The first relation in \eqref{eq:common-scaled} and the
first relation in \eqref{eq:storedreduced} give
\[
          d_i=F(\eta_i+s+t_i)=\eta_iG_i.
\]
The second relation in \eqref{eq:storedreduced} is the
old sum relation. The two common curve relations give
the old unit relation by Lemma~\ref{lem:identity}.
Thus all old equations are recovered. Conversely every
old solution has \(F\ne0\), so
\(x=F^{-1}\) and \(\widehat F=F/c\) satisfy all the new
relations. These choices are unique.

Every monomial in the displayed equations has degree
at most two. No detector or pole equation changes.
We delete \(k\) field coordinates and add two, giving
\(2k-4\) fewer rational coordinates. The old charts
have \(4k\) field equations; the new charts have \(2k+3\).
Taking rational components gives a saving
\(2(4k-(2k+3))=4k-6\).
\end{proof}

The source labels deleted by this theorem are
211 and 212 in every input block. Direct inspection of
the two published source tables shows that they occur
only in \texttt{new1.d\_piF},
\texttt{new1.d\_Gh}, and \texttt{new1.unit}.
Their normalizer and parameter weights vanish, and their
contact weights are one. Thus no hidden substitution
into a detector is needed. The theorem is a local
algebraic replacement; the existence of a shared contact
gauge is a separate arithmetic assertion.

\section{Why the alternate numerator cannot simply be shared}

Let \(P_0\) be a nontorsion point on \(E_D\), and write
\(\Gamma=\Z P_0\). For the common numerator
\(U=[j]P_0\), the alternate points are
\[
                         V_i=U-[n_i]P_0.
\]
All statements below concern finite points with nonzero
abscissa, as required by the existing charts.

\begin{proposition}\label{prop:alternate}
For integer labels \(n_i\):
\begin{enumerate}
\item If \(V_i=V_l\), then \(n_i=n_l\).
\item The function \(y/x\) is injective on the finite
points of \(\Gamma\) with nonzero abscissa. Consequently
sharing the slope \(y(V_i)/x(V_i)\) also forces \(n_i=n_l\).
\item If \(1/x(V_i)=1/x(V_l)\), then either \(n_i=n_l\)
or \(2j=n_i+n_l\). Therefore three distinct labels
cannot all share one reciprocal abscissa.
\end{enumerate}
\end{proposition}
\begin{proof}
The first assertion follows by subtracting the point
equalities and using that \(P_0\) is nontorsion.

For the second put \(T_0=(0,0)\), the nonzero point
of order two. If \(s=y/x\) is fixed, the abscissa
satisfies
\[
                         x^2-s^2x-D=0.
\]
Two points with this slope either have the same
abscissa and are equal, or their abscissas have
product \(-D\). In the latter case the addition
formula gives the other point as
\[
            T_0-Q=\left(-D/x(Q),-Dy(Q)/x(Q)^2\right),
\]
whose slope is again \(s\). Their sum is \(T_0\).
Two elements of the torsion-free group \(\Gamma\)
cannot sum to nonzero torsion. Thus the latter case
does not occur in \(\Gamma\).

Finally equal reciprocal abscissas mean equal
abscissas. On the Weierstrass curve this gives
\(V_l=V_i\) or \(V_l=-V_i\).
The latter equality is equivalent to
\([2j-n_i-n_l]P_0=O\), and hence
\(2j=n_i+n_l\). A fixed abscissa has at most two
points above it, so it cannot contain three distinct
\(V_i\).
\end{proof}

For the quartic input labels
\(n_i=51X_i^2+2504\), inputs \(0,1,2\) give three
distinct labels. This supplies an explicit obstruction
to identifying all their alternate reciprocal
abscissas. These observations do not exclude a
different detector with a different point interpretation.

\section{A restriction-of-scalars obstruction to a one-parameter chart}
\label{sec:restriction}

We now use the actual source curve
\[
 K=\Q(\sqrt2),\qquad d=3-\sqrt2,\qquad
 D=d^2=11-6\sqrt2,\qquad E:y^2=x^3-Dx.
\]
Let \(\overline E\) denote its conjugate under the
nontrivial automorphism of \(K/\Q\). Write
\(\mathcal A=\Res_{K/\Q}E\).
This is an abelian surface with
\(\mathcal A(\Q)=E(K)\).

\begin{lemma}\label{lem:noisogeny}
There is no nonzero \(K\)-homomorphism
\(E\longrightarrow\overline E\).
\end{lemma}
\begin{proof}
Put \(d'=3+\sqrt2\). Since \(dd'=7\), the conjugate
coefficient satisfies
\[
                  \frac{\overline D}{49D}
                      =\left(\frac{d'}7\right)^4
                      =d^{-4}.
\]
Thus \(\overline E\) is \(K\)-isomorphic to the
quadratic 7-twist of \(E\). More explicitly, with
\[
                         u=\frac{\sqrt7}{d},
\]
the map
\[
              \psi:E\longrightarrow\overline E,\qquad
              (x,y)\longmapsto(u^2x,u^3y)
\]
is an isomorphism over \(K(\sqrt7)\), because
\(u^4D=\overline D\).

The geometric endomorphism ring of this \(j=1728\)
elliptic curve is \(\Z[\mathrm i]\), with the
order-four automorphism
\[
                       [\mathrm i](x,y)=(-x,\mathrm i y).
\]
One can see the assertion directly over \(\mathbb C\):
the period lattice is stable under multiplication
by \(\mathrm i\), and hence is homothetic to an ideal
of the principal ideal domain \(\Z[\mathrm i]\).
Its multiplier ring is therefore \(\Z[\mathrm i]\).
Every endomorphism is defined over \(K(\mathrm i)\).

The real quadratic extension \(K(\sqrt7)\) is not
contained in \(K(\mathrm i)\): the maximal real
subfield of the latter is \(K\). Choose an absolute
Galois automorphism \(\sigma\) fixing \(K(\mathrm i)\)
and sending \(\sqrt7\) to \(-\sqrt7\). Then
\(\sigma(\psi)=-\psi\), while \(\sigma\) fixes every
geometric endomorphism of \(E\).

If \(\varphi:E\to\overline E\) were a nonzero
\(K\)-homomorphism, the endomorphism
\(\alpha=\psi^{-1}\varphi\) would satisfy both
\(\sigma(\alpha)=\alpha\) and
\(\sigma(\alpha)=-\alpha\). The torsion-free ring
\(\Z[\mathrm i]\) then gives \(\alpha=0\), a
contradiction.
\end{proof}

\begin{proposition}\label{prop:simple}
The abelian surface \(\mathcal A\) is simple over
\(\Q\).
\end{proposition}
\begin{proof}
If a proper nonzero abelian subvariety \(B\) were
defined over \(\Q\), it would have dimension one.
After base change to \(K\),
\[
                      \mathcal A_K\simeq E\times\overline E.
\]
The two projections of \(B_K\) are conjugate. If one
were zero then both would be zero, contrary to the
inclusion of \(B_K\). Both are therefore nonzero
homomorphisms of elliptic curves, hence isogenies
defined over \(K\). Composing one with the dual of
the other gives a nonzero \(K\)-homomorphism from
\(E\) to \(\overline E\), contradicting
Lemma~\ref{lem:noisogeny}.
\end{proof}

\begin{lemma}\label{lem:nontorsion}
The source point
\[
 P=\left(\frac{57-28\sqrt2}{8},
         \frac{-660+417\sqrt2}{32}\right)
\]
is nontorsion.
\end{lemma}
\begin{proof}
The displayed coordinates satisfy the curve equation
by exact substitution. In \(K_2=\Q_2(\sqrt2)\),
normalize the valuation by \(v(\sqrt2)=1\), so
\(v(2)=2\). The two coordinate valuations are
\[
                         v(x(P))=-6,\qquad v(y(P))=-9.
\]
For the numerator of \(x(P)\), the terms have
valuations \(0\) and \(5\); for the numerator of
\(y(P)\), they have valuations \(4\) and \(1\).
These unequal valuations give the displayed values
after dividing by 8 and 32. Hence
\[
                          v(z(P))=3>v(2).
\]

The integral Weierstrass model has a formal group
at the identity, even though the special fiber
has bad reduction elsewhere. For any integer \(m\),
its multiplication series has integral coefficients
and begins \(mT\). At parameter depth \(v(T)>2\),
the higher terms of \([2](T)\) have valuation at
least \(2v(T)>v(T)+2\), so
\[
                          v([2](T))=v(T)+2.
\]
For odd \(m\), the linear term is a unit times
\(T\) and all higher terms have greater valuation,
so \(v([m](T))=v(T)\). It follows that for every
nonzero integer \(m\),
\[
                  v(z([m]P))=3+2v_2(m)<\infty.
\]
Thus no nonzero multiple of \(P\) is the identity.
\end{proof}

\begin{theorem}[Density of every infinite orbit subset]
\label{thm:dense}
Let \(\widetilde P\in\mathcal A(\Q)\) correspond
to \(P\in E(K)\). Every infinite subset of
\(\Z\widetilde P\) is Zariski dense in
\(\mathcal A\).
\end{theorem}
\begin{proof}
The Zariski closure of the group
\(\Z\widetilde P\) is an algebraic subgroup defined
over \(\Q\). Its identity component has positive
dimension by Lemma~\ref{lem:nontorsion}, so
Proposition~\ref{prop:simple} makes it all of
\(\mathcal A\). The same is true for the closure
of \(m\Z\widetilde P\) whenever \(m\ne0\).

Apply the classical Mordell--Lang theorem for
abelian varieties to a closed subvariety \(V\)
and the finitely generated group
\(\Gamma=\Z\widetilde P\). In the form used here,
\(V\cap\Gamma\) is a finite union of cosets of
subgroups of \(\Gamma\), each contained in \(V\)
\cite{Faltings,Mazur}. If \(V\cap\Gamma\) is
infinite, at least one such subgroup is nonzero.
Every nonzero subgroup of \(\Gamma\simeq\Z\) is
\(m\Gamma\) for some \(m\ge1\), whose Zariski
closure is all of \(\mathcal A\).
Its coset therefore forces \(V=\mathcal A\).
\end{proof}

\begin{corollary}[Precisely scoped chart obstruction]
\label{cor:oneparameter}
No input-independent finite algebraic chart over
one rational base parameter, nor any fixed finite
collection of such charts, can contain all the
points
\[
                        [(x^2+1)]P,\qquad x\in\Z .
\]
The same statement holds for any infinite set of
distinct integer multiples of \(P\).
\end{corollary}
\begin{proof}
By a chart of the indicated kind we mean a finite
algebra over an open subset of \(\mathbb A^1_\Q\),
with finitely many algebraic generators, optional
nonvanishing guards, and an algebraic map on the
guarded open subset to \(\mathcal A\).
A finite algebra has the same dimension as its
base, at most one. The closure of its image
therefore has dimension at most one. Finitely
many charts give a finite union of such proper
closed subsets of the two-dimensional
\(\mathcal A\).

The points in the displayed set are infinite
and distinct after replacing \(x\) by nonnegative
integers. Theorem~\ref{thm:dense} makes that set
Zariski dense in \(\mathcal A\), which is
incompatible with the preceding image bound.
Nonreduced finite algebras and extra finite root
stages do not alter the dimension argument.
\end{proof}

\paragraph{Scope of the obstruction.}
The qualifier \emph{input-independent} is essential.
If the chart coefficients also depend on the existing
input \(X\), its base may be
\(\operatorname{Spec}\Q[X,t]\), of dimension two;
the argument above then gives no obstruction.
The same caveat applies to dependence on another
variable shared by the whole construction. The
integer indices used to choose elliptic multiples
are not automatically rational base coordinates:
one cannot count an integer label as an
algebraically available variable without an
actual polynomial construction. Nothing here
is a lower bound for the number of variables
in an arbitrary Diophantine definition of
\(\Z\) or in an undecidability reduction.

\section{A local check against an unsafe detector replacement}

One tempting way to remove the old contact weight
of \(G\) would replace its \(c^2G^2\) contribution
by \(-7c^4G^2\). A nonsquare condition on 7 at
the newly chosen completeness primes does not
justify the same change in soundness. The consumed
primes in the soundness argument range outside
a finite exceptional set; they are not all
required to make 7 nonsquare.

Here is a concrete failure of the attempted
noncancellation rule. At \(p=113\), choose an
embedding \(K\hookrightarrow\Q_{113}\) and
formal curve points with parameters
\[
                 z(V)=25p,\qquad z(H)=p^2,\qquad c=p^{-1}.
\]
The formal expansion is
\(f(Q)=z(Q)^2+O(z(Q)^6)\). Consequently
\[
 G=f(V)=25^2p^2+O(p^6),\qquad
 f(H)=p^4+O(p^{12}),
\]
and
\[
        1+c^4f(H)^4-7c^4G^2
                  =1-7\cdot25^4+O(p^4).
\]
Since
\[
       1-7\cdot25^4=-2734374,\qquad
       v_{113}(2734374)=1,
\]
this expression has positive valuation one.
Thus the old residue-class argument that each
factor has nonpositive valuation cannot be
reused from reciprocal parity alone.
This is a counterexample to that proposed
local noncancellation inference, not a
counterexample to every possible redesigned
complete system. A repair would need an
additional soundness argument; no improved
numerical bound is claimed from this replacement.

\section{Reproducibility and current obligations}

The exact symbolic identities, the twist ratio,
the displayed point's curve equation and local
valuations, and the \(113\)-adic integer check
are replayed by
\nolinkurl{undecidability/check_point_chart_lemmas.py}.
The geometric endomorphism and Mordell--Lang
arguments are mathematical proofs and are not
asserted to be certified by that computation.

The algebraic saving of
Theorem~\ref{thm:stored} is available when the
common-contact-gauge theorem and its full
quadratic compilation have been supplied.
The simple alternate-point and one-parameter
obstructions show why several more direct
identifications cannot complete the task.
These results preserve the distinction between
proved chart transformations, arithmetic
interface hypotheses, and unresolved attempts.

\begin{thebibliography}{9}
\bibitem{physical}
\(\mathbb Z\mathbb Q\) Diophantine Project,
\emph{A quartic in 193 rational coordinates},
v74, 17 November 2026.
\nolinkurl{papers/quartic/quartic193_20261117_v74.tex}.

\bibitem{scalar}
\(\mathbb Z\mathbb Q\) Diophantine Project,
\emph{A scalar-point sharing construction in 163
and 179 variables}, v82, 20 November 2026.
\nolinkurl{papers/quartic/scalar_point163_179_20261120_v82.tex}.
Complete maps:
\nolinkurl{computation/variable_reductions_v82/general/}.

\bibitem{Faltings}
G. Faltings, \emph{The general case of S. Lang's
conjecture}, in V. Cristante and W. Messing (eds.),
\emph{Barsotti Symposium in Algebraic Geometry},
Perspectives in Mathematics 15, Academic Press,
1994, pp.\ 175--182.
\url{https://doi.org/10.1016/B978-0-12-197270-7.50012-7}.

\bibitem{Mazur}
B. Mazur, \emph{Abelian varieties and the
Mordell--Lang conjecture}, in
\emph{Model Theory, Algebra, and Geometry},
MSRI Publications 39, 2000, pp.\ 199--227,
Theorem 3.1 and the discussion following it.
\url{https://library2.msri.org/books/Book39/files/mazur.pdf}.
The explicit characteristic-zero formulation
was also read at the university-hosted mirror
\url{https://math.unm.edu/~buium/Mazurmordelllang.pdf}.
\end{thebibliography}
\end{document}
